%!TEX program = xelatex
\documentclass{xmu}
%%%%%%%%%%%%%%%%%%%%%%%%%%%%%%%%%%%%%%%%%%%%%%%%%%%%%%%%%%
%%%%%%%%    XMU Undergraduates Thesis Template    %%%%%%%%
%%%%%%%%             Made by: F5Soft              %%%%%%%%
%%%%%%%%  https://github.com/F5Soft/xmu-template  %%%%%%%%
%%%%%%%%%%%%%%%%%%%%%%%%%%%%%%%%%%%%%%%%%%%%%%%%%%%%%%%%%%

\begin{document}

% 电子版 / 打印版（取消注释下一行即为打印版）
% \print
% 取消注释后，将在某些偶数页产生空白页，使得下一部分的内容从奇数页开始

% 取消注释后，仅使用数字作为章的编号
% \arabicchapter

% 毕业设计 / 毕业论文（取消注释下一行即为毕业设计）
\design

% 主修 / 辅修（取消注释下一行即为辅修）
% \minor

% 标题
\title{基于RAG技术的智能标书撰写系统}
{Intelligent Bid Document Writing System Based on RAG Technology}

% 姓名
\author{朱纪玉}

% 学号
\idn{22920212204333}

% 学院
\college{信息学院}

% 专业
\subject{网络空间安全}

% 年级
\grade{2021级}

% 校内指导教师
\teacher{罗晔\; 副教授}

% 校外指导教师（注释则不显示）
% \otherteacher{F5Soft\; 教授}

% 完成时间
\pubdate{二〇二五年四月三十日}

% 关键词
\keywords{检索增强生成；大语言模型；标书}
{Undergraduate Thesis; LaTeX; Xiamen University}

% 生成封面、承诺书
\maketitle

% 如果想要将致谢放在最前，可将最后的致谢移动至此

%%%%%%%% 摘要 %%%%%%%%

% 中文摘要
\begin{abstract}
啥也没写
\end{abstract}

% 英文摘要
\begin{enabstract}
empty
\end{enabstract}

% 生成中英文目录
\tableofcontents

%%%%%%%% 正文 %%%%%%%%
\xmuchapter{绪论}{Introduction}
\xmusection{招投标行业数字化发展现状与挑战}{}
% #region
随着我国经济社会的快速发展，招投标行业作为市场经济的重要组成部分，其作用日益凸显。
招投标行业的健康发展对于促进经济增长、提高市场竞争力、维护社会公平正义具有重要意义。
招投标行业具有广泛的市场覆盖面，包括基础设施建设、公共服务、工业生产等多个领域。
据统计，2024年，我国招标采购规模约34.74万亿元，其中，政府公开采购招标规模2.49万亿元，工程招标采购26.82万亿元，其他采购招标5.43万亿元；\cite{chyxx}
目前，在国家进一步“放管服”、加速推进“优化营商环境”的大背景下，招标采购加速向信息化、数字化发展。\cite{zhaotoubiaoxianzhuang}
一方面，招标采购交易全流程逐渐实现网络信息化、数字化:2024年全国电子招投标平台日均处理量突破100万单，AI评标系统在中部六省试点准确率达92\%\cite{report}。
另一方面，招标投标的服务模式从单一、程序性服务向复合型、综合性咨询服务转型。
\par
\begin{figure}[!htb]
    \centering
    \includegraphics[width=30em]{专项公共服务.png}\\
    \caption{中国招标投标公共服务平台已经支持CA互认和远程异地评标和评标专家库共享}\label{xmulogo}
\end{figure}
目前，地方政府公共资源交易平台功能日臻完善，各大央企集团大多已建成使用招标采购电子交易平台，行业领域内招标采购业务活动基本实现全流程在线化，
目前正在加快向数字化和智能化方向转型，AI应用基础良好。
总体而言，与招标人相比，投标人对人工智能技术应用驱动力不足，在功能开发和应用方面均不如招标人及资源交易中心广泛和深入。
投标人也已经逐步意识到大语言模型对招标采购行业的深远影响，目前已经有投标人使用GPT编制投标文件部分章节，如施工组织方案。\cite{AIzhanwang}
\par
%#endregion
% #region
从开发者的角度来看，开发为投标人服务的智能工具存在以下难点：
1.招标、投标文件往往包含大量的专业术语，需要遵循特定的逻辑或行业规则，目前的自然语言处理技术理解和生成这些内容时存在一定局限性。
2.在投标过程中，投标人需要综合市场、企业、客户、竞争对手等多方面信息进行决策，仅仅依靠自然语言处理技术或简单的检索技术难以提供有参考价值的决策建议。
3.招标采购行业涉及大量的商业机密和敏感信息，尤其是对于提供SAAS（Software as a Service，软件即服务）的服务商而言，在技术上和社会信任上都需要长期努力。
% #endregion
\xmusection{RAG技术}{}
% #region
\xmusubsection{RAG简介}{}
RAG（Retrieval-Augmented Generation，检索增强生成）
是一种连接外部数据源以增强大语言模型（LLM）输出质量的技术。
这种技术帮助 LLM 访问私有数据或特定领域的数据，并解决幻觉问题。
因此，RAG 已被广泛用于许多通用的生成式 AI（GenAI）应用中，
如 AI 聊天机器人和推荐系统。
RAG技术一般可以划分为信息检索和生成两个阶段。
\par
在信息检索阶段，系统解析用户输入的文本，
从构建的专业知识库中，通过向量检索、图检索等方式筛选相关程度最高的文本片段。
在生成阶段，模型将检索到的文档片段作为上下文与用户输入的文本一起输入到seq2seq模型中，得到输出。
\par
\begin{figure}[!htb]
    \centering
    \includegraphics[width=30em]{RAG.png}\\
    \caption{RAG技术基本原理}\label{xmulogo}
\end{figure}
\xmusubsection{RAG技术在智能标书撰写工具中的应用潜力}{}
标书撰写的过程可以分解为招标要求理解和标书内容生成两个任务，
传统的大语言模型能够满足一般的文本理解和内容生成任务，但受限于模型的训练数据，
对于标书撰写这样涉及专业知识和企业信息以及多种综合因素的复杂任务，存在一定局限性。
而RAG技术通过检索和利用外部知识库中的相关信息，
能够为标书撰写提供更丰富的知识支持和更准确的生成内容。

RAG技术预计可以实现以下功能：
1、理解招标要求：通过检索与招标文件相关的行业标准、法规和历史案例，
RAG模型能够更精准地理解招标文件的核心要求，并利用相关知识指导内容生成。
2、生成标书内容：结合检索到的外部知识和输入的招标文件要求，RAG模型可以生成可靠的标书正文内容，
并且依靠检索提供的证据链，为人工复核提供便利。
3、提供决策参考：RAG模型可以检索市场数据、竞争对手信息等外部知识，结合预设的竞价、博弈模型，
为标书撰写提供有价值的决策参考。
% #endregion
\xmusection{本文工作}{}
% #region
本文工作主要包括以下几个方面：
分析现有的智能标书撰写工具，设计基于RAG技术的智能标书撰写系统；
实现“招标需求理解-标书框架生成-标书内容生成”工作流；
实现基于知识图谱的方案库，将用户上传的文件，解析为可用于RAG的知识图谱；
实现基于Prompt的风格迁移功能，为用户提供预设或自定义的风格选项；
实现“数据检索-工具调用-输出决策”的决策参考工作流；
% #endregion
\xmusection{论文结构}{}
% #region
第一章：
本章介绍招投标行业的数字化发展现状，阐述为投标人服务的智能工具的开发难点。
随后，简要介绍RAG技术并分析标书撰写中的应用潜力。
接着，概述本文的主要工作内容，和论文的整体结构。
\par
第二章：
本章介绍传统的招投标方法和现有的智能标书撰写工具，分析其优缺点。
详细回顾RAG技术的相关研究进展，其中重点介绍基于知识图谱的RAG技术。
\par
第三章：
本章将对智能标书撰写系统进行需求分析。随后阐述系统架构设计。
\par
第四章：
本章将描述系统的开发环境和工具，详细介绍系统功能的具体实现。
\par
第五章：
本章将总结本文的主要研究成果，并讨论系统的优势和可能的改进方向。
% #endregion
\xmuchapter{相关工作}{}
\xmusection{招投标领域的研究现状}{}
\xmusubsection{传统投标方法}{}
% #region
一般标书的结构分为商务部分、技术部分和报价部分（但招标文件特殊要求格式除外）。

商务部分：一般包括投标人说明、厂家介绍、业绩、合同、产品授权书、法人授权书、三证、资格证书、交货期、付款方式、售后服务、承诺书、商务偏离表、商务应答、备品备件专用工具清单等，要严格按照标书内容要求及顺序编写。

技术部分：一般包括投标设备技术说明、图纸设计、技术参数、产品配置、技术规格偏离表，技术力量简介、安装施工方案、产品质量、产品简介、产品彩页等等，要严格按照标书内容要求及顺序编写。
其中技术偏离表需要标注正偏离、负偏离、无偏离；如投标产品的技术指标优于招标要求即为正偏离，反之为负偏离，符合招标要求即为无偏离；

报价部分，一般包括报价一览表、分项报价表、投标函、投标承诺书、投标单位名称、投标代表签字、法人代表签字等，要严格按照招标文件要求及顺序编写。
\par
标书编写可以分为初步目录编排和后期目录编排两种方式。
初次目录编排：根据招标文件的要求，初步编写投标文件目录；为了方便收集投标书的资料。对评分点、控标点、优势应在初步目录中标注，其目的为了让标书制作者重视该部分文档。
后期目录编排：按事先拟定好的投标文件目录，对正文内容的标题设置为标题1、标题2、标题3、标题4等，然后自动生成投标文件的目录。
\par
在标书电子版制作完成后，参与人员需要进行交叉检查，对错误的地方进行指正修改；需要重点审查的地方有：
正本制作正副本内容一致；
字体、格式是否统一；
审查报价产品名细是否符合招标产品需求名细（包括产品型号和数量），分项（分包）报价是否符合和正确；
审查报价表中的分报价和总报价的计算、大小写是否正确，报价表、投标一览表、投标函中的报价大小写是否一致；
审查开标文件书写格式是否与招标方的要求一致；
\par
在标书打印完成后，需签字处完成签字盖章，最终审查无误后，进行密封。
\par
传统的标书编写方法依赖人工操作，为确保标书的质量和合规性，投标团队需要投入大量的时间和精力进行信息收集、内容撰写和格式调整，存在较多重复性、繁琐性的任务。
开发智能化的标书撰写工具，将投标工作中的重复性和繁琐性任务自动化，有望提高标书编写的效率和准确性。
% 随着招投标行业的数字化转型，越来越多的企业开始寻求智能化的解决方案，以提高标书编写的效率和质量。
% #endregion
\xmusubsection{现有智能标书撰写工具}{}
% #region
目前，市面上已经出现了一批智能标书撰写工具，如链企标书、星火投标等。
这些工具的开发背景各不相同，部分由传统招标采购行业的专业公司开发，部分由互联网科技巨头或新兴的互联网创业公司开发。
在服务形式方面，
大多数工具提供SaaS（Software as a Service，软件即服务）模式，
用户通过网络浏览器登录平台使用标书撰写工具，
少数工具提供本地部署服务。

这些智能标书撰写工具的核心功能基本相同，仅在功能完成度和用户体验上存在差异。
总体而言，这些工具具备以下几项核心功能：
招标文件解析：解析用户上传的投标文件，提取标书撰写的参考信息。
投标目录生成：根据自动提取或用户输入的参考信息，生成标书目录。
投标正文生成：根据标书目录，生成标书内容。
除此以外，部分工具有在线编辑、风格迁移、文件云端存储引用、封面模板等附加功能。

\begin{figure}[!htb]
    \centering
    \includegraphics[width=38em]{现有工具.png}\\
    \caption{现有AI工具的核心功能，图中以链企AI为例}\label{xmulogo}
\end{figure}

目前市面上的智能标书撰写工具在实际应用中仍然存在如下问题。

1、生成的内容不够专业。
标书撰写涉及复杂的细分领域，
不同行业的招标要求和标书撰写规范存在较大差异。
现有工具仅仅能区分服务类、货物类、工程类、其它类四大类别。

2、无法提供决策参考。
标书中包含必不可少的报价部分，现有工具几乎不能给出有可靠依据的报价。

% #endregion
\xmusubsection{RAG技术研究现状}{}
% #region
RAG技术最早由Facebook AI Research提出，
随后在任务规划、推理检索、多模态等研究方向进一步发展。

Rewrite-Retrieve-Read模型利用 LLM 的能力，
通过重写模块和 LM 反馈机制来优化检索查询，并通过更新重写模型来提高任务性能\cite{RRR}。
HyDE通过关注生成答案与真实文档之间的嵌入相似性，旨在提高检索的相关性\cite{HyDE}。
2024年7月2日，微软开源GraphRAG\cite{graphRAG}。
GraphRAG利用基于知识图谱的方法，从原文档导出知识图谱，聚类后生成社区摘要。
将社区摘要对用户问题的响应汇总，作为对用户的最终响应。
这一方法在生成答案的全面性、多样性都相较于传统RAG的基线有显著提升。

经过多年发展，RAG技术的范式已经从朴素RAG发展到模块化RAG，
RAG过程已经被抽象出各种专用子组件，
如搜索模块、记忆模块、预测模块、任务适配器模块。\cite{RAG}

RAG技术栈也越加丰富。LangChain和LLamalndex等大模型框架已经支持RAG技术的快速实现。
目前也已经出现多个成熟的专用RAG框架，如Haystack、RAGFlow。

% #endregion
\xmuchapter{智能标书撰写系统设计}{}
\xmusection{需求分析}{}
\xmusection{系统架构设计}{}
% #region

系统架构主要分为三个层次：

1.数据访问层：负责存储和管理企业的历史标书、招标文档、用户参考文段等数据资源。这一层包括方案库（知识图谱）、资历库和素材库（半结构化数据），为企业信息管理和标书内容生成提供数据支持。

2.业务逻辑层：是系统的核心，负责标书内容的生成和处理。这一层包括RAG（Retrieval-Augmented Generation）模块、标书内容生成模块（招标理解、目录生成、正文生成）、风格迁移模块（基于prompt和基于LoRA）以及后期处理模块。

3.表示层：提供用户界面，使用户能够方便地输入招标文件、管理模板、查看生成的标书内容等。

系统包含以下模块
1.企业信息管理模块：该模块负责存储管理企业的历史标书、方案库（知识图谱）和资历库素材库（半结构化数据），提供可视化管理页面。

2.标书内容生成模块：该模块是系统的核心，负责根据招标文件的要求生成标书的框架和正文内容，包括招标理解、目录生成和正文生成三个子模块。

3.风格迁移模块：该模块实现基于prompt的风格迁移技术，提供预定义的风格和用户自定义风格。

4.决策参考模块：该模块提供成本定价模型、盈亏平衡定价和第一价格密封拍卖算法等决策，为用户提供决策参考。

5.后期处理模块：该模块对生成的标书内容进行校对、格式调整和优化，以文件的格式输出。

% #endregion
\xmusection{企业信息管理}{}

\xmusection{方案库（知识图谱）}{}
% #region
基于知识图谱的方案库，作为将用户上传的文件解析为可以用于RAG的结构化数据的模块，
需要实现文件解析和检索支持两方面的功能。
我们需要从文件中提取两部分内容，
首先将文本原文分块并向量化存储，向量将作为全局索引用于检索，
其次从文本中提取实体和关系的语义信息
这两部分内容都将存储在知识图谱中。

文本解析工作流如图：

\begin{figure}[!htb]
    \centering
    \includegraphics[width=38em]{pipline.png}\\
    \caption{上传的文件解析为文本后，分块进行实体和关系提取，合并相同实体后存储为知识图谱}\label{}
\end{figure}

文档解析：将用户上传的文件解析为文本。\par
文本分块：过长的文本会降低文本中早期信息的召回率\cite{Long}，
因此需要将文件提取的文本进一步分块。\par
实体和关系提取：利用LLM提取文本块中的实体和关系信息。\par
文本向量化：将文本块语义信息转化为向量信息，以便于后续检索。\par
知识图谱构建：将文本块、实体作为节点，关系作为边存储在知识图谱中。
实体通过“FROM\_CHUNK”关系连接到所属的文本块。
文本块节点之间通过“NEXT\_CHUNK”关系作为边连接。
同一文件中不同文本块提取的实体和关系直接视为同一实体进行合并。

检索工作流如图：

\begin{figure}[!htb]
    \centering
    \includegraphics[width=38em]{search.png}\\
    \caption{将用户输入转化为查询语句，先通过向量检索得到初始文本块，再搜寻向量节点形成完整图}\label{}
\end{figure}

语句转化：将用户上输入的自然语言语句转化为知识图谱查询语句。\par
向量检索：从全局索引（即向量信息）中检索找到初始文本块。\par
图检索：根据初始文本块在相应的节点中查询邻居节点路径。\par
返回结果：将检索到的图输出。

知识图谱参数定义如下:

\begin{table}[!htb]
    \centering
    \caption{节点定义}
    \label{Node}
    \begin{tabular}{|l|l|l|l|l|}
        \hline
        \bf\songti 参数 & \bf\songti 参数类型& \bf\songti 参数描述 & \bf\songti 备注 \\ \hline
        id             & str         &                 &             \\ \hline
        label               & str          &             &               \\ \hline
        properties               & dict[str,Any]         &                 &               \\ \hline
        embedding\_properties             & dict[str, list[float]] | None       &       &               \\ \hline
    \end{tabular}
\end{table}

\begin{table}[!htb]
    \centering
    \caption{关系定义}
    \label{Relationship}
    \begin{tabular}{|l|l|l|l|l|}
        \hline
        \bf\songti 参数 & \bf\songti 参数类型& \bf\songti 参数描述 & \bf\songti 备注 \\ \hline
        start\_node\_id             & str         &                 &             \\ \hline
        end\_node\_id               & str          &             &               \\ \hline
        type               & str          &             &               \\ \hline
        properties               & dict[str,Any]         &                 &               \\ \hline
        embedding\_properties             & dict[str, list[float]] | None       &       &               \\ \hline
    \end{tabular}
\end{table}


% #endregion
\xmusection{资历库和素材库}{}
% #region
资历库存储企业的合同、授权书、三证等资质信息。素材库存储其它撰写标书涉及的信息。
资历库需要能够存储和检索企业的资质信息。
一般来说，资历的文件格式为PDF或图片，仅依靠文件名无法有效检索，
因此需要将文件的信息提取并关联存储。

各个行业涉及的证书信息不同，需要灵活的存储方式。
将资历库设计为半结构化的数据库。
对于常见的资质文件，可以预设存储的字段。
对于不在预设中的资质文件，可以由用户自行增加存储字段。

资历库存储以JSON格式，结构如下：

\begin{verbatim}
{
    "id": "string",
    "name": "string",
    "type": "string",
    "file\_path": "string",
    "properties": {
        "key1": "value1",
        "key2": "value2"
    }
}
\end{verbatim}

% #endregion
\xmusection{标书内容生成}{}
\xmusubsection{识别招标需求}{}
% #region
该模块负责提取招标文件中的关键信息，用于指导后续的标书撰写。
该模块分为两步骤：文件解析和文本解析。
一般而言，招标文件为PDF或word文档，在少数情况下为图片格式（纸质文档拍照）。
文件解析需要将文件转化为文本格式，文本解析需要将文本分块并提取关键信息。

文件解析过程在后端处理。文本解析由LLM完成。

为了方便处理，我们将提取到的信息分为“项目概况”和“评分标准”两类
（实际的招标文件中可能涉及更复杂的关键信息）。
项目概况包括招标单位、招标项目、招标范围、投标人资格要求、投标文件递交时间和地点、开标时间和地点等信息。
评分标准包括商务评分、技术评分、报价评分等信息。

% #endregion
\xmusubsection{生成标书目录}{}
% #region
该模块负责生成标书的目录。根据“项目概况”和“评分标准”两类信息，生成标书目录（尚未引入RAG）。
实际生成过程中，用户可能需要调整目录的顺序和内容，为保证灵活性，我们采用逐章节生成的方式。
具体生成方案如下：

初始章节生成：LLM仅根据全局信息（即“项目概况”和“评分标准”和“风格信息”）生成第一章目录。
后续章节生成：将全局信息、已经生成的全部章名、上一章完整目录、下一章完整目录（如果存在）作为输入，生成当前章的目录。

为方便处理，我们将标书的框架固定为“章-节-小节”三级标题形式。每个章可对应多个节，每个节可对应多个小节。
在不进行任何生成时下，用户可以在前端对目录进行任意调整修改。
% #endregion
\xmusubsection{生成标书正文}{}
% #region
该模块负责生成标书正文。每个小节根据全局信息、当前“章-节-小节”三级标题进行基于RAG的生成。
初始生成时，仅通过方案库进行RAG检索。LLM的输出为“内容”“引用列表”固定格式。

用户可以对生成的内容进行如下修改：
编辑：用户可以直接在前端编辑生成的内容。
重写生成：用户可以选择某部分内容重新生成。全局信息、当前“章-节-小节”三级标题、重写部分上下文将作为输入。
引用生成：用户可以手动添加来自方案库、素材库、资历库的引用，手动添加引用生成将不再进行RAG检索。
报价生成：用户可以通过“决策参考”模块生成报价部分的内容。
风格化生成：用户可以覆盖全局的风格化设置，选择当前小节的风格化选项。

% #endregion
\xmusection{风格迁移}{}
% #region
该模块负责调整生成标书风格。
风格迁移采用基于Prompt的方式，将预设文字或用户自定义的文字作为提示词，
指导LLM生成符合特定风格的内容。
% #endregion
\xmusection{决策参考}{}
% #region

% #endregion
\xmusection{后期处理}{}
% #region

% #endregion
\xmuchapter{系统实现与实验}{}
    \xmusection{系统开发环境与工具}{}
        \xmusubsection{系统环境}{}
        \begin{itemize}
            \item 操作系统：Ubuntu 22.04
            \item Python版本：3.8
            \item 依赖库：deepseek、fastapi、uvicorn、pydantic、pymongo、neo4j
        \end{itemize}
        \xmusubsection{开发工具}{}
        \begin{itemize}
            \item IDE：VSCode
            \item 版本控制：Git
            \item 文档编写：Markdown
        \end{itemize}
        \xmusubsection{系统框架}{}
        \begin{itemize}
            \item 前端框架：Vue.js
            \item 后端框架：FastAPI
            \item 数据存储：MongoDB、neo4j
        \end{itemize}
    \xmusection{系统功能实现}{}
        \xmusubsection{方案库(知识图谱)}{}
        使用neo4j存储知识图谱。neo4j是一个图数据库，使用cipher语句进行数据存储和查询。\par
        后端使用neo4j-graphrag库和neo4j库进行知识图谱操作。
        neo4j库用于连接到neo4j进行基本数据操作。
        neo4j-graphrag库是neo4j提供的知识图谱框架库，用于实现文件解析和检索支持。
        文件解析过程中，文本向量化使用阿里百炼text-embedding-v3模型，向量维度为1024维，
        实体和关系提取使用deepseek-chat模型。

        前端使用阿里g6图可视化引擎。
        该引擎提供了图的绘制、布局、分析、交互、动画等基础的图可视化能力。

        后端api定义：
        \begin{itemize}
            \item POST /api/upload：
        \end{itemize}

        测试样例：(待修改)
        "The son of Duke Leto Atreides and the Lady Jessica, Paul is the heir of House "
        "Atreides, an aristocratic family that rules the planet Caladan."
        \begin{figure}[!htb]
            \centering
            \includegraphics[width=38em]{grapheg.png}\\
            \caption{文段生成的知识图谱，0节点为文本块}\label{}
        \end{figure}

        \xmusubsection{招标需求识别模块}{}
        % 定义招标需求的结构化数据模型，使用deepseek api进行招标需求识别。
        % 定义Chapter、Section、Subsection数据模型，使用pydantic进行数据验证。
        % 利用langchain结构化输出模型，将大模型输出转化为结构化数据。
        % 后端api使用fastapi实现，前端使用axios进行数据交互。
        % api定义：
        % \begin{itemize}
        %     \item POST /api/identify：识别招标需求
        %     \item GET /api/chapters：获取章节列表
        %     \item GET /api/sections：获取节列表
        %     \item GET /api/subsections：获取小节列表
        % \end{itemize}
        \xmusubsection{标书目录生成模块}{}
    %     使用deepseek-chat模型生成目录。
    %     \begin{itemize}
    %         \item 目录生成：根据招标需求和评分标准，生成“章-节-小节”固定格式的目录。
    %         \item 目录保存：将生成的目录保存为结构化数据，以便后续使用。
    %         \item 目录审核：用户对生成的目录进行审核和修改，并提供反馈信息。
    %     \end{itemize}
    %     \xmusubsection{标书正文生成模块}{}
    %     使用方案库后端作检索增强，使用deepseek-chat模型生成正文。
    %     \begin{itemize}
    %         \item 目录生成：根据招标需求和评分标准，生成“章-节-小节”固定格式的目录。
    %         \item 目录保存：将生成的目录保存为结构化数据，以便后续使用。
    %         \item 目录审核：用户对生成的目录进行审核和修改，并提供反馈信息。
    %     \end{itemize}
    %     \xmusubsection{风格迁移模块}{}
    %     在生成正文时，根据用户提供的示例文本或描述性提示，选择合适的风格迁移方法（基于prompt或基于LoRA），对生成的标书内容进行风格迁移。
    %     api定义：
    %     \begin{itemize}
    %         \item POST /api/style-transfer：风格迁移
    %         \item GET /api/styles：获取可用的风格列表
    %     \end{itemize}
    %     \xmusubsection{决策参考模块}{}
    %     在标书生成报价部分，使用成本定价模型、盈亏平衡定价和第一价格密封拍卖算法等决策参考功能，辅助用户进行投标决策。
    %     \xmusubsection{后期处理模块}{}
    %     对生成的标书内容进行校对、格式调整和优化，确保标书的准确性和专业性。
    %     \xmusubsection{模板管理模块}{}
    %     支持用户自定义和保存标书模板，以便于未来的标书撰写。
    % \xmusection{系统测试}{}
    %     \xmusubsection{测试方法}{}
    %     \xmusubsection{测试结果}{}
\xmuchapter{总结与展望}{Conclusion and Future Work}
未来优化点：使用graphRAG的全局-社区-摘要模式
尝试多种文本块节点的连接和存储方式
知识图谱融合
识别招标需求可以更细致

资历库素材库没做
决策参考实现
% #region
% \xmuchapter{使用说明}{Introduction}
% \xmusection{环境配置}{Environment Setup}
% 首先，需要配置好LaTeX环境。这里推荐TeX Live + VSCode。
% 其中TeX Live可访问{https://www.tug.org/texlive/}下载；
% VSCode安装完成后需要再搜索安装LaTeX Workshop插件。
% \xmusection{编译}{Compiling}
% 环境配置完成后，可以从GitHub上下载本模版的代码\cite{template}。
% 下载完成后，可尝试打开example.tex文件，点击右上角的绿色运行按钮进行编译。
% 在macOS系统上，编译一般会失败，这是因为使用的pdflatex找不到对应中文字体的结果。
% 可以进入左侧的LaTeX Workshop面板，选择Recipe为latexmk (xelatex)，点击运行，即可编译成功。
% \xmusection{使用模版}{Use Templates}
% 本模版根据《厦门大学本科毕业论文（设计）规范》\cite{xmuthesis}封装了大量格式设置等内容，
% 使用时无需考虑各种格式指令，只需设置好章节标题，填充摘要、附录、参考文献、致谢等内容即可。
% \xmusubsection{示例论文}{Example Paper}
% example.tex文件内有详细的注释和使用示例，可以直接在该示例基础上修改。
% 论文中的所有插图需统一放在figures文件夹内，并根据文件名来包含。
% \xmusubsection{自定义模版}{Customization}
% xmu.cls文件是模版文件，如果您具有排版强迫症，想和我一样对模版的细节进行微调，
% 可以自己修改xmu.cls模版文件。模版文件中也提供了大量注释便于您进行修改。
% 如果您认为修改的效果不错，可以在GitHub上\cite{template}发起Pull request贡献代码。
% 您也可以发起Issue提出要添加的功能。
% \par
% 修改模版需要有一定的LaTeX基础，可以参考《LaTeX 2e完全学习手册》
% \cite{latex2e}和《简单粗暴LaTeX》\cite{easylatex}中的内容进行修改。
% \xmuchapter{使用示例}{Examples}
% \xmusection{二级标题}{Section (English Ver.)}
% 一级标题（章）总会另起一页。
% \xmusubsection{三级标题}{Subsection (English Ver.)}
% 通过xmuchapter、xmusection、xmusubsection指令指定一级（章）、二级（节）、三级（小节）标题。
% 因为英文目录的需要，这三个指令均接收两个参数，
% 第一个参数为中文的标题内容，第二个参数为英文的标题内容。
% \subsubsection{四级标题}
% 如果需要四级标题的话，不会出现在目录内。
% \xmusection{图表}{Figures and Tables}
% \xmusubsection{插图}{Figures}
% 图片编号均按照章节号-图片号进行编号。
% 可通过caption指令设置图片小标题，
% 可通过label和ref指令引用图片编号。
% 厦门大学校徽如下图 \ref{xmulogo} 所示。
% \begin{figure}[!htb]
%     \centering
%     \includegraphics[width=12em]{xmu-logo-icon.jpg}\\
%     \caption{厦门大学校徽}\label{xmulogo}
% \end{figure}
% 根据《厦门大学本科毕业论文（设计）规范》\cite{xmuthesis}要求，
% 图片的标题和图号需要在图的下方正中位置。
% 下面再展示一个两图并排的示例：
% \begin{figure}[!htb]
%     \begin{minipage}{0.5\linewidth}
%         \centering
%         \includegraphics[height=5em]{xmu-logo-icon.jpg}
%         \caption{厦门大学校徽}\label{xmu1}
%     \end{minipage}
%     \begin{minipage}{0.5\linewidth}
%         \centering
%         \includegraphics[height=5em]{xmu-logo-text.jpg}
%         \caption{厦门大学文字}\label{xmu2}
%     \end{minipage}
% \end{figure}
% \xmusubsection{表格}{Tables}
% 对于表格，编号也按照章节号-表格号进行编号。
% 同样可通过caption指令设置表格小标题，可通过label和ref指令引用表格编号。
% 下表 \ref{database} 中展示了一个数据库的表设计作为示例。
% \begin{table}[!htb]
%     \centering
%     \caption{基本资料表}
%     \label{database}
%     \begin{tabular}{|l|l|l|l|l|}
%         \hline
%         \bf\songti 字段名称 & \bf\songti 字段类型 & \bf\songti 长度 & \bf\songti 字段描述 & \bf\songti 备注 \\ \hline
%         账户号             & Number          & 30            &                 & 主键            \\ \hline
%         密码              & Number          & 30            & 加密              &               \\ \hline
%         姓名              & Varchar         & 50            &                 &               \\ \hline
%         电子邮箱            & Varchar         & 50            & VIP客户必填         &               \\ \hline
%     \end{tabular}
% \end{table}
% 根据《厦门大学本科毕业论文（设计）规范》\cite{xmuthesis}要求，
% 表格的标题和图号需要在表的上方正中位置。
% 表格内容的字体为5号宋体，表头的字体为5号宋体加粗。
% \xmusection{画图}{Plotting}
% 可以使用tikzpicture环境绘制函数、折线图、散点图等，例如图 \ref{plot} 所示：
% \begin{figure}[!htbp]
%     \centering
%     \begin{tikzpicture}
%         \begin{axis}[
%                 xlabel={Temperature [Celsius Degree]},
%                 ylabel={Solubility [g per 100 g water]},
%                 xmin=0, xmax=100,
%                 ymin=0, ymax=120,
%                 xtick={0,20,40,60,80,100},
%                 ytick={0,20,40,60,80,100,120},
%                 legend pos=north west,
%                 ymajorgrids=true,
%                 grid style=dashed,
%             ]

%             \addplot[
%                 color=blue,
%                 mark=square,
%             ]
%             coordinates {
%                     (0,23.1)(10,27.5)(20,32)(30,37.8)(40,44.6)(60,61.8)(80,83.8)(100,114)
%                 };
%             \legend{CuSO\(_4\cdot\)5H\(_2\)O}

%         \end{axis}
%     \end{tikzpicture}
%     \caption{Temperature dependence of CuSO\(_4\cdot\)5H\(_2\)O solubility}\label{plot}
% \end{figure}
% \xmusection{公式}{Formulas}
% 公式可以分为行内公式 $a^2+b^2=c^2$ 和行间公式：
% $$
%     \int_a^b\int_a^x\int_x^{x^2}(3y+7)z\;\mathrm{d}z\,\mathrm{d}y\,\mathrm{d}x
% $$
% $$
%     |-0.75|^{2}+(5.5-4.2)^2+\sqrt[4.2]{3}+\log_{3.3}{4.7}+\cos1.24\,\sin1.7\approx5.17
% $$
% 如果需要指定公式编号，同样可通过label和ref指令来引用。公式 \ref{eq1} 的正确结果应该为 $\frac{\sqrt{3}}{6}$。
% \begin{equation}\label{eq1}
%     \lim_{x\to0}\frac{\displaystyle{\int_{\sin x}^x\sqrt{3+t^2}\,\mathrm{d}t}}{x(e^{x^2}-1)}
% \end{equation}
% 公式 \ref{eq2} 提供了一个矩阵示例，注意加粗字母最好是使用boldsymbol而非bm：
% \begin{equation}\label{eq2}
%     \boldsymbol{Ax} = \boldsymbol{b},\quad\boldsymbol{A} = \left(\begin{array}{ccc}
%             a_{11} & a_{12} & a_{13} \\
%             a_{21} & a_{22} & a_{23} \\
%             a_{31} & a_{32} & a_{33} \\
%         \end{array}\right),\quad \boldsymbol{b}=\left(\begin{array}{c}
%             b_1 \\ b_2 \\ b_3\end{array}\right)
% \end{equation}
% 公式 \ref{eq3} 提供了一个行列式示例：
% \begin{equation}\label{eq3}
%     \left|\begin{array}{cc}
%         1 & 2 \\
%         3 & 4 \\
%     \end{array}\right|=1\times4-2\times3=-2
% \end{equation}
% 公式 \ref{eq4} 提供了一个方程组示例：
% \begin{equation}\label{eq4}
%     \begin{cases}
%         mgh=mv_1^2                         \\
%         \dfrac12mv_2^2-\dfrac12mv_1^2=F_fs \\
%         F_f=\mu mg                         \\
%     \end{cases} \Rightarrow s,v_1,F_f
% \end{equation}
% 公式字体采用的是mathptmx。
% \xmusection{算法}{Algorithms}
% 如果只是展示代码，可考虑使用verbatim环境
% \footnote{更高级的代码展示可以考虑使用listings宏包中的lstlisting环境，
% 支持代码高亮、换行、标注行号、自定义颜色、自定义字体等功能。}。
% 这里插入了一个脚注。可以通过footnote指令插入脚注\footnote{脚注示例。}。
% \begin{verbatim}
% #include <unistd.h>
% int main()
% {
%     write(STDOUT_FILENO, "Hello, world!\n", 14);
%     return 0;
% }
% \end{verbatim}
% 如果只是展示算法，而不展示具体代码，可以使用算法环境实现，例如算法 \ref{alg1}：
% \begin{algorithm}[!htb]
%     \KwData{sum, $i$: integer}
%     \Begin{
%         \For{$i\leftarrow0$ \KwTo $100$ \do}{
%             sum $\leftarrow$ sum + $i$\;
%             \If{\rm sum $>4000$}{
%                 \bf break\;
%             }
%         }
%         \bf print \rm$i$\;
%     }
%     \caption{循环求和算法}\label{alg1}
% \end{algorithm}
% #endregion
%%%%%%%% 参考文献 %%%%%%%%

\begin{reference}
    % 如果需要使用 bib 文件导入参考文献，则取消注释下一行
    % \bibliography{references.bib}
    % bib 文件可以通过百度学术、Google Scholar 的引用界面自动生成
    % 已自动按照 GB/T 7714-2005 设置参考文献的引用格式

    % 手动添加参考文献
    \begin{thebibliography}{5}
        % \bibitem[1]{xmuischool} 厦门大学图书馆. i学堂直播：使用LaTeX排版论文全攻略-进阶篇[EB/OL]. \url{https://library.xmu.edu.cn/wd/jzkj/i_xt.htm}, 2021.
        % \bibitem[2]{xmuthesis} 厦门大学教务处. 厦门大学本科毕业论文（设计）规范[Z]. 厦门大学教务处: 厦大教〔2016〕5号.
        % \bibitem[3]{latex2e} 胡伟. LaTeX 2e完全学习手册[M]. 清华大学出版社, 2013.
        % \bibitem[4]{easylatex} 吴康隆. 简单高效LaTeX[M]. 人民邮电出版社, 2020.
        % \bibitem[5]{template} F5Soft.site. 厦门大学本科毕业论文LaTeX模版[EB/OL]. \url{https://github.com/F5Soft/xmu-template}, 2022.
        \bibitem[1]{zhaotoubiaoxianzhuang}芮丽梅,石国虎.新形势下招标投标现状与发展趋势[J].招标采购管理,2021,(07):11-12.
        \bibitem[2]{chyxx}智研咨询.2025-2031年中国招标采购行业发展模式分析及未来前景规划报告[EB/OL]. \url{https://www.chyxx.com/research/202110/980866.html}, 2024.
        \bibitem[3]{report}中研普华产业研究院.2025-2030年招投标市场深度分析及发展趋势研究咨询报告[EB/OL]. \url{https://www.chinairn.com/news/20250220/103457473.shtml}, 2024.
        \bibitem[4]{AIzhanwang}程建宁.AI技术在招标投标活动中的应用及展望[J].招标采购管理,2023,(07):61-63.
        \bibitem[5]{RAG}Yunfan Gao, Yun Xiong, Xinyu Gao, Kangxiang Jia, Jinliu Pan, Yuxi Bi, Yi Dai, Jiawei Sun, Meng Wang, Haofen Wang.Retrieval-Augmented Generation for Large Language Models: A Survey[J/OL].arXiv:2312.10997[cs.CL]
        \bibitem[6]{graphRAG}Darren Edge, Ha Trinh, Newman Cheng, Joshua Bradley, Alex Chao, Apurva Mody, Steven Truitt, Dasha Metropolitansky, Robert Osazuwa Ness, Jonathan Larson.From Local to Global: A Graph RAG Approach to Query-Focused Summarization[J/OL].arXiv:2404.16130[cs.CL]
        \bibitem[7]{RRR}Xinbei Ma, Yeyun Gong, Pengcheng He, Hai Zhao, Nan Duan.Query Rewriting for Retrieval-Augmented Large Language Models[J/OL].arXiv:2305.14283 [cs.CL]
        \bibitem[8]{HyDE}Luyu Gao, Xueguang Ma, Jimmy Lin, Jamie Callan.Precise Zero-Shot Dense Retrieval without Relevance Labels[J/OL]arXiv:2212.10496 [cs.IR]
        \bibitem[9]{Long}Yuri Kuratov, Aydar Bulatov, Petr Anokhin, Dmitry Sorokin, Artyom Sorokin, Mikhail Burtsev.In Search of Needles in a 11M Haystack: Recurrent Memory Finds What LLMs Miss[J/OL].arXiv:2402.10790 [cs.CL]
    \end{thebibliography}
\end{reference}

%%%%%%%% 附录 %%%%%%%%

\begin{appendix}
    % \xmusection{附表}{Tables}
    % 这里放一些附录的内容，例如表格或其他说明。如附表\ref{atable}。
    % \begin{table}[!htb]
    %     \centering
    %     \caption{附表示例}
    %     \label{atable}
    %     \begin{tabular}{llr}
    %         \hline
    %         \multicolumn{2}{c}{Item} &                          \\ \cline{1-2}
    %         Animal                   & Description & Price (\$) \\ \hline
    %         Gnat                     & per gram    & 13.65      \\
    %                                  & each        & 0.01       \\
    %         Gnu                      & stuffed     & 92.50      \\
    %         Emu                      & stuffed     & 33.33      \\
    %         Armadillo                & frozen      & 8.99       \\ \hline
    %     \end{tabular}
    % \end{table}
    % 附录中的编号均以字母A开头。例如公式 \ref{eq5}
    % \begin{equation}\label{eq5}
    %     \sin \left(2x+\frac\pi2\right)=\frac 12
    % \end{equation}
    % \xmusection{厦门大学本科毕业论文（设计）规范摘要}{Abstract of Xiamen University Undergraduate Dissertation (Design) Specification}
    % 主修专业毕业论文（设计）封面使用 160g 白色双胶纸，辅修封面为 160g 浅黄色皮纹纸。内页均为 A4 规格 80g 双胶纸。\par
    % 章的标题占2行，标题以外的文字为1.5倍行距。\par
    % 上边距和左边距应留 25mm 以上间隙，下边距和右边距应分别留 20mm 以上间隙。\par
    % 每页须加“页眉”和“页码”。奇数页页眉内容为当前章名，如“第一章 绪论”。偶数页页眉内容为论文题目。学位论文的页码，正文、参考文献、附录部分用阿拉伯数字连续编码并居中，前置部分用罗马数字单独连续编码居中（封面除外）。\par
    % 封面中文标题：二号黑体\par
    % 封面英文标题：三号 Times New Roman 加粗\par
    % 中文摘要标题：小三号黑体\par
    % 中文关键词标题：小四号黑体\par
    % 中文摘要、关键词内容：小四号宋体\par
    % 英文摘要标题：小三号 Times New Roman 加粗\par
    % 英文关键词标题：小四号 Times New Roman 加粗\par
    % 英文摘要、关键词内容：小四号 Times New Roman\par
    % 中文目录标题：小三号黑体\par
    % 中文目录中章的标题：四号黑体\par
    % 中文目录中节的标题：小四号黑体\par
    % 中文目录中三级标题：小四号宋体\par
    % 英文目录标题：小三号 Times New Roman 加粗\par
    % 英文目录中章的标题：四号 Times New Roman 加粗\par
    % 英文目录中节的标题：小四号 Times New Roman 加粗\par
    % 英文目录中三级标题：小四号 Times New Roman\par
    % 章的标题：小三号黑体\par
    % 节的标题：四号黑体\par
    % 三级标题：小四号黑体\par
    % 正文：小四号宋体\par
    % 页眉：小五号宋体\par
    % 页码：小五号 Times New Roman\par
    % 注释内容：小五号宋体\par
    % 表格、图的标题、单位、表头：五号宋体加粗\par
    % 表格内容：五号宋体\par
    % 表格、图的资料来源：小五号宋体\par
    % 参考文献标题：小三号黑体\par
    % 中文参考文献表：五号宋体\par
    % 英文参考文献表：五号 Times New Roman\par
    % 附录标题：小三号黑体\par
    % 致谢标题：小三号黑体\par
    % 致谢内容：小四号宋体\par
    % % 对于中英文混杂的内容，{\bfseries\songti 中文的字体若是用宋体，英文的字体则采用Times New Roman}；{\sffamily 中文的字体若是黑体，英文的字体则采用Arial}。\footnote{因此，本模版提供了如下的几种字体设置。rm：罗马字族，中文宋体，英文Times。sf：等线字族，中文黑体，英文Arial。bf：粗体序列，中文黑体，英文Times加粗。it：斜体序列，中文楷体，英文Times斜体。如果需要使用难看的宋体加粗，则需要通过 \textbackslash bf\textbackslash songti实现。当然，也可以使用 songti、heiti、kaishu 手动指定中文字体，此时英文字体均默认为Times。}
    % \xmusubsection{附录三级标题}{Subsection in Appendix}
    % \subsubsection{附录四级标题}
\end{appendix}

%%%%%%%% 致谢（默认在最后） %%%%%%%%

\begin{acknowledgement}
暂时不写
\end{acknowledgement}

\end{document}
